\documentclass[9pt,a4paper,twocolumn]{ltjsarticle}
\usepackage[top=20mm,bottom=22mm,left=17mm,right=17mm,columnsep=8mm]{geometry}
\usepackage{amsmath}
\usepackage[hidelinks]{hyperref}
\newcommand{\SchemaCommit}{c2b4de1}
\newcommand{\NSamples}{442}
\newcommand{\NFeatures}{10}
\newcommand{\NTrain}{353}
\newcommand{\NTest}{89}
\newcommand{\NFullRuns}{1}
\newcommand{\TestMse}{3424.26}
\newcommand{\CpuCount}{4}
\newcommand{\MemoryGib}{16}
\newcommand{\RepositoryCommit}{156ce2c}
\expandafter\def\csname lineC00\endcsname{47}
\expandafter\def\csname lineC01\endcsname{73}
\expandafter\def\csname lineC02\endcsname{100}
\expandafter\def\csname lineC03\endcsname{119}
\expandafter\def\csname lineC04\endcsname{134}
\expandafter\def\csname lineC05\endcsname{153}
\expandafter\def\csname lineC06\endcsname{172}
\expandafter\def\csname lineC07\endcsname{191}
\expandafter\def\csname lineC08\endcsname{210}
\expandafter\def\csname lineC09\endcsname{229}
\expandafter\def\csname lineC10\endcsname{244}
\expandafter\def\csname lineR01\endcsname{259}
\expandafter\def\csname lineR02\endcsname{292}
\expandafter\def\csname lineR03\endcsname{308}
\expandafter\def\csname lineR04\endcsname{324}
\expandafter\def\csname lineR05\endcsname{340}
\expandafter\def\csname lineR06\endcsname{397}
\expandafter\def\csname lineR07\endcsname{415}
\expandafter\def\csname lineR08\endcsname{433}
\expandafter\def\csname lineR09\endcsname{458}
\expandafter\def\csname lineR10\endcsname{476}
\expandafter\def\csname lineR11\endcsname{494}
\expandafter\def\csname lineR12\endcsname{516}
\expandafter\def\csname lineR13\endcsname{534}
\expandafter\def\csname lineR14\endcsname{549}
\expandafter\def\csname lineR15\endcsname{567}
\expandafter\def\csname lineR16\endcsname{582}
\expandafter\def\csname lineR17\endcsname{600}
\expandafter\def\csname lineR18\endcsname{615}
\expandafter\def\csname lineR19\endcsname{630}
\expandafter\def\csname lineR20\endcsname{648}
\expandafter\def\csname lineR21\endcsname{666}
\expandafter\def\csname lineR22\endcsname{681}
\expandafter\def\csname lineoutputTestMse\endcsname{701}
\expandafter\def\csname lineobservedTrainIndex\endcsname{863}
\expandafter\def\csname lineobservedTestIndex\endcsname{1231}
\expandafter\def\csname linefullRunTestMses\endcsname{1648}

% 各文の元になった claims の ID。スキーマ（claims.json、検証の結果を含む commit）の、その claim の行へリンクする
\newcommand{\schemaurl}{https://github.com/auto-res2/formal-first-diabetes-regression-4/blob/\SchemaCommit/formal/claims.json}
\newcommand{\cid}[1]{\textsuperscript{\textsf{\scriptsize[\href{\schemaurl\#L\csname line#1\endcsname}{#1}]}}}
% 検証された値。スキーマの、その値の行へリンクする
\newcommand{\val}[1]{\textsuperscript{\textsf{\scriptsize[\href{\schemaurl\#L\csname line#1\endcsname}{値}]}}}
\newcommand{\code}[1]{\texttt{\small #1}}
\begin{document}

\twocolumn[{
\begin{center}
{\Large\bfseries 糖尿病データセットに対する線形回帰のテスト MSE と\\しきい値 3400 の比較}\\[2mm]
{\normalsize auto-res2}
\end{center}
\vspace{3mm}
}]

\section{主張}

実験の前に，次の主張を宣言した．\cid{C01}Efron らの糖尿病データセットに線形回帰モデルを当てはめる実験が出力する，テストデータに対する MSE は，3400 を下回る．
\cid{C00}あわせて，この主張が成り立ち，かつ，信頼性条件（\ref{sec:reliability} 節）を満たすことを主張した．
検証の結果，出力された MSE は 3400 を下回らず，信頼性条件は満たされた（\ref{sec:result} 節）．

\section{実験設定}

\cid{C02}MSE は，次の宣言の手順で，1 つの値に定まる．宣言は，実験の前に置いた．

\paragraph{データセット}
\cid{C03}Efron ら~\cite{efron} の糖尿病データセットを，配布元の URL\footnote{\url{https://www4.stat.ncsu.edu/~boos/var.select/diabetes.tab.txt}}から取得して使用する．取得するファイルの SHA-256 を固定した．データは \NSamples{} 件で，説明変数は \NFeatures{} 個（AGE，SEX，BMI，BP，S1〜S6），目的変数は Y である．

\paragraph{読み込みと分割}
\cid{C04}ファイルは \code{numpy.loadtxt} で読み込み（1 行目の列名は読み飛ばす），1〜10 列目を説明変数，11 列目を目的変数とする．読み込んだ値は加工しない．
\cid{C05}\code{sklearn.model\_selection.train\_test\_split} を，乱数のシード 0，テストデータの割合 20\% で使用し，テストデータ以外の行を訓練データとする．

\paragraph{当てはめ，予測，評価}
\cid{C06}訓練データの説明変数 \NFeatures{} 個すべてと目的変数に，\code{LinearRegression.fit} を既定のパラメータで適用する．
\cid{C07}当てはめたモデルと説明変数に \code{LinearRegression.predict} を適用して，予測値を求める．
\cid{C08}テストデータの目的変数と予測値に \code{sklearn.metrics.mean\_squared\_error} を適用した値を，MSE とする．

\paragraph{計算環境とリポジトリの版}
\cid{C09}GitHub Actions の標準の runner（ubuntu-24.04，x86\_64，\CpuCount{} コア，メモリ \MemoryGib{} GiB）で実行する．
\cid{C10}実験は，リポジトリ\footnote{\url{https://github.com/auto-res2/formal-first-diabetes-regression-4}}の commit \code{\RepositoryCommit} を取り出して実行する．実験コード，ライブラリの版を固定するファイル，Python の版を含む実行のコマンドは，すべてこの commit にある．

\section{結果}
\label{sec:result}

実験は \NFullRuns{} 回だけ実行した．出力された MSE は \TestMse\val{outputTestMse} であった．
\cid{C01}この値は 3400 を下回らないので，主張は偽である．この判定は値の比較だけで決まり，公理に依存しない．
\cid{C00}したがって，あわせた主張も偽である．偽になったのは MSE が 3400 を下回らなかったためで，信頼性条件は次節のとおり満たされた．

\section{信頼性条件}
\label{sec:reliability}

\cid{R01}信頼性条件は，下の 3 つの公理，宣言と観測の一致，HARKing をしていないこと，チェリーピッキングをしていないことからなる．検証の結果，すべて成り立った．この判定は，3 つの公理に依存する．

\paragraph{宣言と観測の一致}
\cid{R05}すべての宣言について，宣言した値は，GitHub の検証器が観測した値と一致し，その宣言のほかの検証も成り立った．
\cid{R06}データセット（URL，ファイルの SHA-256，件数，列）は，宣言と一致した．\cid{R07}宣言した URL から取得したファイルの SHA-256 も，宣言と一致した．\cid{R08}したがって，下の公理のもとで，実行が読んだファイルは Efron らのデータセットである．
\cid{R09}読み込みの関数と引数は，宣言と一致した．
\cid{R10}分割の関数と引数は，宣言と一致した．\cid{R11}訓練データ \NTrain{} 件\val{observedTrainIndex}とテストデータ \NTest{} 件\val{observedTestIndex}は重ならず，合わせると全 \NSamples{} 件になる．
\cid{R12}当てはめの関数，パラメータ，説明変数の数は，宣言と一致した．\cid{R13}当てはめに渡したデータは，読み込んだファイルの訓練データの行そのものである．
\cid{R14}予測の関数は，宣言と一致した．\cid{R15}予測に渡した説明変数は，読み込んだファイルのテストデータの行そのものである．
\cid{R16}MSE の計算の関数は，宣言と一致した．\cid{R17}渡した目的変数は，読み込んだファイルのテストデータの行そのものである．\cid{R18}渡した予測値は，予測の関数が返した値そのものである．
\cid{R19}計算環境（runner の種類，アーキテクチャ，コア数，メモリ）は，宣言と一致した．
\cid{R20}実行が取り出したリポジトリの commit は，宣言と一致した．

\paragraph{公理}
次の 3 つは，検証器が無いため，正しいものとして置いた．
\cid{R02}Efron らが用いた糖尿病データセットは，配布元の URL から取得したファイルである．
\cid{R03}実験コードは，仕様を満たす．すなわち，使ったものと計算した MSE をそのまま出力に記録し，宣言のとおりにデータを関数へ渡す．
\cid{R04}ライブラリは信用できる．すなわち，宣言のとおりに使えば，ライブラリが計算する MSE は，宣言したとおりに計算した MSE である．関数が何を計算するかは主張しない．

宣言と観測の一致は，出力に記録された値についての検証である．記録が実際の実行を表すことは，実験コードの公理による．

\paragraph{研究公正}
\cid{R21}本番の実行は，失敗したものも含めて，すべて，判定に使った宣言と同じ宣言の commit から起動した．宣言は，実験の後に書き換えていない．
\cid{R22}成功した本番の実行は \NFullRuns{} 回\val{fullRunTestMses}で，出力された MSE は，判定に使った値と一致する．

\begin{thebibliography}{9}
\small
\bibitem{efron} B. Efron, T. Hastie, I. Johnstone, R. Tibshirani: Least angle regression, Annals of Statistics, 32(2), 407--499, 2004.
\end{thebibliography}

\end{document}
